\documentclass[11pt]{article}
\usepackage{fontspec}
\usepackage[french]{babel}
\usepackage{amsmath}
\usepackage{amssymb}
\usepackage{amsthm}
\usepackage{mathrsfs}
\usepackage{tikz-cd}
\usepackage{enumitem}

\newcommand{\printedpage}[1]{\marginpar{\footnotesize p.~#1}}
\newcommand{\msp}[1]{\marginpar{\footnotesize [#1]}}
\newcommand{\limind}{\varinjlim}
\newcommand{\limproj}{\varprojlim}
\newcommand{\Fl}{\mathrm{Fl}}
\newcommand{\Is}{\mathrm{Is}}
\newcommand{\Ho}{\mathrm{Ho}}
\newcommand{\qpair}{\mathrel{\leftarrow\mkern-4mu\mapstochar\mkern-2mu\rightarrow}}

\begin{document}

\printedpage{24}
% head: A. Grothendieck : Dérivateurs, Ch. XIX, Édition des Universités de Montpellier II et Paris VII
% foot: Version 11 février 2019
\msp{page 43}
instructif, bien sûr, si $\mathcal{M}$ est stable par petites $\limind$ , que ses objets soient accessibles etc. et que $W$ soit L-accessible (i.e. stable par petites $\limind$ filtrantes, et de plus accessible \dots). Le premier candidat pour $I$ auquel je pense serait un $B_G$ ($G$ un groupe, p.ex. $G=\pm1$), c'est en effet l'exemple aux antipodes des ensembles ordonnés, lesquels, eux, ne semblent pas faire trop de problème (moyennant l'hypothèse noethérienne tout au moins).

Les exemples suivants ne sont pas aussi convaincants que ça, hélas!

\textbf{Exemple 3.10.} Soit $\mathcal{M}$ une catégorie quelconque, alors les couples $(\Fl(\mathcal{M}),\Is(\mathcal{M}))$ et $(\Is(\mathcal{M}),\Fl(\mathcal{M}))$ sont tous deux des couples de Quillen,
\[
\Fl(\mathcal{M}) \qpair \Is(\mathcal{M}) \ , \qquad \Is(\mathcal{M}) \qpair \Fl(\mathcal{M}) \ . \tag*{(3.10.1)} % ?? uncertain: shape of the arrow -- left-right arrow with a vertical stroke in the middle (as printed), exact glyph of the edition unknown
\]
Il s'ensuit que si on pose
\[
W \ = \ \Is(\mathcal{M}) \ , \qquad C \ = \ F \ = \ \Fl(\mathcal{M}) \ , \tag*{(3.10.2)}
\]
$(W,C,F)$ est un triple de Quillen clos \textsuperscript{(}\footnote[22]{au moins si on admet que $\mathcal{M}$ est stable par $\limind$ et $\limproj$ \textit{finies.}}\textsuperscript{)}. D'ailleurs pour toute $I$ dans Cat, on a $W(I)=\Is(\mathcal{M}(I))$, donc
\[
\Ho_W(I) \ = \ \mathcal{M}(I) \ . \tag*{(3.10.3)}
\]
Donc la question de l'existence des foncteurs $u_!^W$, $u_*^W$ se réduit à celle des foncteurs $u_!$, $u_*$ habituels, adjoints à gauche et à droite de $u^* : \mathcal{M}(I') \longrightarrow \mathcal{M}(I)$.

\msp{page 44}
On sait bien que l'existence de ces foncteurs pour $u$ donné, fait appel à l'existence de certaines $\limind$ resp. $\limproj$ dans $\mathcal{M}$; p.ex. si $I'=e$, il est nécessaire et suffisant que les $\limind$ (resp. $\limproj$ ) de type $I^\circ$ soient représentables dans $\mathcal{M}$. Mais avec les hypothèses faites sur $\mathcal{M}$, il n'y a pas de problème si $u : I \longrightarrow I'$ est un foncteur entre catégories \textit{finies.} On va donner un exemple où il n'en est plus ainsi.

\textbf{Exemple 3.11.} Soit $k$ un anneau commutatif, et $\mathcal{M}=K^{\mathrm{b}}(k\text{-Ab})$ la catégorie des complexes de cochaînes bornés de $k$-modules, $W$ l'ensemble des quasi-isomorphismes, donc
\[
\Ho_W \ = \ D^{\mathrm{b}}(k\text{-Ab}) \ ,
\]
catégorie dérivée bornée de la catégorie abélienne des $k$-modules. On sait, en utilisant l'existence d'assez de projectifs dans $k$-Ab, que cette structure est associée à une structure de Quillen close $(W,C,F)$, décrite par \textsc{Quillen}
\msp{page 45}
dans Homotopical Algebra -- à cela près qu'il prend $D^-(k)$ au lieu de $D^{\mathrm{b}}(k)$. Peut-être, pour vérifier les axiomes, doit-on supposer que $k$ est de dimension cohomologique finie.

Prenons $I=B_G$, où $G$ est un groupe (p.ex. $G=\{\pm1\}$), et $u : I \longrightarrow e$. Alors $\mathcal{M}(I)$ est la catégorie des complexes bornés de $k[G]$-modules. On voit que $u_!^W$ resp. $u_*^W$ existe si et seulement si pour tout tel complexe borné, ou simplement pour un $k[G]$-module sans plus $M$, le complexe $L\Gamma^G(M)$ resp. $R\Gamma_G(M)$ appartient à $D^{\mathrm{b}}(k)$, ou ce qui revient au même,
\printedpage{25}
% head: A. Grothendieck : Dérivateurs, Ch. XIX, Édition des Universités de Montpellier II et Paris VII
% foot: Version 11 février 2019
si et seulement si on a
\[
\begin{aligned}
H_i(G,M) \ &= \ 0 \quad \text{pour } i \gg 0 \ ,\\
\text{resp.}\quad H^i(G,M) \ &= \ 0 \quad \text{pour } i \gg 0 \ ,
\end{aligned}
\]
pour tout $k[G]$-module $M$. Ce n'est pratiquement jamais le cas; quand on prend p.ex. $k=\mathbf{Z}$ et $G$ groupe fini, cela n'est le cas que si $G=e$. L'exemple le plus économique me paraît $G=\pm1$, $k=\mathbf{F}_2$.

\msp{page 46}
Donc ici $u_!^W$ et $u_*^W$ n'existent pas, pour $I \longrightarrow e$ avec $I$ catégorie finie. Mais aussi $\mathcal{M}$ n'est pas stable par petites $\limind$ .

\textbf{Exemple 3.12.} Soit $\mathcal{M}$ une catégorie muni d'un localiseur $W$ fortement saturé, tel que $\mathcal{M} \longrightarrow \underbrace{\mathcal{M}W^{-1}}_{=\mathcal{H}}$ admette un adjoint -- à gauche disons. On a donc une situation de deux foncteurs adjoints % [sic] muni -> munie, agreement with « une catégorie »
\[
\begin{tikzcd}[column sep=huge]
\mathcal{H} \arrow[r, shift left, "v_!\ \text{pleinement fidèle}"] & \mathcal{M} \arrow[l, shift left, "v^*\ \text{localisation}"]
\end{tikzcd}
\ , \tag*{(3.12.1)}
\]
et $W=(v^*)^{-1}(\mathrm{Iso}(\mathcal{H}))$. On trouve que pour toute $I$ dans Cat, on a
\[
\mathrm{HO}_W(I) \ \simeq \ \mathcal{H}(I) \ \overset{\text{déf}}{=} \ \mathrm{Hom}(I^\circ,\mathcal{H}). \tag*{(3.12.2)} % [sic] HO_W -> Ho_W, notation Ho_W used in (3.10.3) and Exemple 3.11
\]
On peut donc songer à faire un exemple où dans $\mathcal{M}$ les petites $\limind$ et $\limproj$ existent, mais non dans $\mathcal{H}$, de sorte que les $u_!^W$ ou $u_*^W$ n'existent pas toujours. Mais ce semble difficile à construire, vu que dans la situation envisagé, on voit facilement que les petites $\limproj$ existent dans $\mathcal{H}$. Donc si dans $\mathcal{H}$ % [sic] envisagé -> envisagée, agreement with « la situation »
\msp{page 47}
il existe une petite famille cogénératrice (et pour ceci, il suffit qu'une telle famille existe dans $\mathcal{M}$), il y aura également des petites $\limind$ , en vertu de SGA 4 I 8.12.7. De même, elles existent si $\mathcal{H}$ est une sous-catégorie pleine de $\mathcal{M}$ \textit{stable par petites} $\limind$ , et il paraît pratiquement impossible de faire un exemple (3.12.1) sans que $\mathcal{H}$ ait cette propriété de stabilité. Sans compter qu'il y a même la perplexité suivante :

\marginpar{! \textbar}\textbf{Question 3.13.} \textit{Existe-t-il une $\mathfrak{U}$-catégorie (\,$\mathfrak{U}$ un univers donné) où les petites $\limind$ soient représentables, sans qu'il en soit de même des petites $\limproj$ ?} \textbf{N.B.} Une telle catégorie n'admet pas de petite famille génératrice, a fortiori elle n'est pas équivalente à une petite catégorie.

\msp{page 48}
\textbf{Question 3.14.} \textsc{Anderson} prétend que pour le prédérivateur défini par une catégorie à cofibrations satisfaisant l'axiome Cof$_4$ transfini, les foncteurs $u_!^W$ existent. Sa démonstration est totalement canulée, néanmoins je n'arrive pas à faire de contre-exemple. Le résultat qu'il annonce serait-il vrai? Je ne connais pas même d'exemple d'un localiseur $W$, dans une catégorie $\mathcal{M}$ \textit{avec petites $\limind$ et petites $\limproj$ ,} tel que $u_!^W$ ou $u_*^W$ n'existent pas toujours -- mais je suis persuadé qu'il y a des exemples (même avec $\mathcal{M}$ une petite catégorie?). J'en suis moins convaincu si on suppose $W$ stable par $\limind$ filtrantes, et au besoin, de plus, que les éléments de $\mathcal{M}$ sont accessibles (plus éventuellement d'autres propriétés catégoriques anodines sur $\mathcal{M}$, telles l'existence d'une filtration
\printedpage{26}
% head: A. Grothendieck : Dérivateurs, Ch. XIX, Édition des Universités de Montpellier II et Paris VII
% foot: Version 11 février 2019
\msp{page 49}
cardinale etc.)

\textbf{Question 3.15.} Cela me ramène à la question particulièrement élusive et irritante, de l'existence des petits ensembles Q-générateurs, dans une partie $\Phi$ de $\Fl(\mathcal{M})$, quand on suppose (disons) que $\mathcal{M}$ est stable par petites $\limproj$ , que ses éléments sont accessibles et qu'il admet une filtration cardinale. Ainsi on a la question : une partie L-accessible de $\Fl(\mathcal{M})$, qu'on suppose au besoin Q-saturé à gauche (ou ce qui revient au même, contenant les isomorphismes, et stable par cochangement de base, par composition et par facteurs directs) admet-elle une \textit{petite partie} $\Phi_0$ Q-génératrice à gauche? % [sic] Q-saturé -> Q-saturée, agreement with « une partie »

\marginpar{! \textbar}D'autre part, je n'ai pas à présent construit un seul exemple d'un couple $\Phi \longleftrightarrow \Psi$, dans une catégorie $\mathcal{M}$, qui ne soit factorisant! Le cas où $\mathcal{M}$ est spéciale, était bien sûr le plus intéressant. Visiblement, dans tout
\msp{page 50}
ça ce sont les exemples qui manquent. Le clef d'une compréhension semble bien être dans les questions de factorisation entre ensembles $\Phi \longleftrightarrow \Psi$. Il faudrait en faire une étude systématique, mais je recule devant la perspective d'une digression qui me prendrait peut-être encore des semaines! % [sic] Le clef -> La clef, « clef » is feminine

\end{document}
